% !TEX root = ../Abschlussarbeit_Jan_Höndl.tex
\section{Einleitung}
\subsection{Problemstellung} \label{sec:problemstellung}
Die Baubranche stellt eine der tragenden Säulen der deutschen Wirtschaft dar, sieht sich jedoch
im Vergleich zu anderen Industriezweigen mit einer signifikanten Digitalisierungslücke konfrontiert \cite{BBSR}.
Während in der Planungsphase durch Methoden wie Building Information Modeling (BIM) bereits digitale
Fortschritte erzielt werden, hinkt insbesondere die Ausführungsphase auf der Baustelle deutlich hinterher.
Dieser Rückstand führt in der Praxis zu einer Vielzahl von Ineffizienzen, die den Projekterfolg bei
Großbauvorhaben gefährden.
\newline
Ein zentrales Hemmnis ist die mangelnde Interkonnektivität der eingesetzten Softwaresysteme.
Da einheitliche Standards und Schnittstellen oft fehlen, entstehen isolierte Datensilos.
Dies führt regelmäßig zum sogenannten „Excel-Fallback“ \cite[Kap.~4.1]{feth2023studie}:
Um den Informationsfluss zwischen den unterschiedlichen Systemen und verschiedenen Akteuren aufrechtzuerhalten,
werden Daten manuell in Tabellenkalkulationen übertragen. Dieser Prozess ist nicht nur fehleranfällig,
sondern verhindert auch eine revisionssichere Dokumentation und Echtzeit-Auswertungen.

Zusätzlich verschärft das Problem des Vendor-Locks die Situation \cite[Kap.~1]{feth2023studie}.
Unternehmen sind oft an geschlossene Ökosysteme einzelner Softwareanbieter gebunden,
was die Erweiterbarkeit um spezialisierte Funktionen oder die Integration innovativer Drittanbieter-Lösungen massiv erschwert.
Bestehende Anwendungen agieren häufig als Insellösungen, die den spezifischen Anforderungen komplexer Bauprozesse, wie einem durchgängigen Inbetriebnahmemanagement, nicht gerecht werden.
Das Inbetriebnahmemanagement stellt während der Bauphase und bei der Inbetriebnahme die Qualitätssicherung dar. Ziel ist es, die vertraglich vereinbarten Leistungen, sowie die behördlich festgelegten Anforderungen nachzuweisen
um eine rechtssichere, geordnete Übergabe an den Bauherren zu gewährleisten. Da dieser Prozess eine Vielzahl unterschiedlicher Gewerke mit meist unterschiedlichen Softwarelösungen involviert \cite[Kap.~4.1]{feth2023studie},
blockiert die fehlende Kompatibilität den notwendigen, systemübergreifenden Informationsaustausch.
\subsection{Forschungslücke}
Trotz der Verfügbarkeit diverser Softwarelösungen im BIM-Kontext, primär in den Bereichen des digitalen Bautagebuchs und der Mängelerfassung, weisen diese oftmals Defizite hinsichtlich ihrer Flexibilität und Adaptierbarkeit an spezifische
Projektanforderungen auf. Zur Lösung dieser Problematik ist eine modulare Architektur erforderlich, die es ermöglicht, spezialisierte Anwendungen flexibel als Add-ons zu integrieren, sodass diese nahtlos auf bestehende Daten zugreifen können.
Um dies zu ermöglichen, bedarf es einer standardisierten und transparenten Schnittstelle, um einen problemlosen Informationsaustausch zu gewährleisten.
Es existieren bereits Forschungen und etablierte Ansätze für graphbasierte Domänenmodelle sowie die Nutzung von GraphQL als Vermittlungsschicht, um die komplexen Strukturen des IFC-Standards abzubilden. Arbeiten wie \cite{wolf:2025:ifc2graphql} zeigen bereits,
dass die Überführung von IFC-Daten in ein GraphQL-Schema grundlegend machbar ist, jedoch eine exakte 1:1-Abbildung an ihre Grenzen hinsichtlich der Komplexität stößt. Trotzdem wird deutlich, dass eine solche Überführung Vorteile in Form von Effizienz
und Entwicklerfreundlichkeit mit sich bringt.

Es fehlt jedoch an einer Verknüpfung dieser Modelle mit einer modularen Architektur als Frontend, mit welcher die beschriebenen Probleme aus \ref{sec:problemstellung} gezielt angegangen werden können.
\subsection{Forschungsfragen}
Um die oben genannten Zielsetzungen methodisch zu untersuchen, konzentriert sich diese Arbeit auf die Beantwortung der folgenden zwei Kernfragen:
\begin{description}
    \item[FF1:] Wie lässt sich eine Graph-basierte Datenstruktur auf Basis des IFC-Standards gestalten, um die komplexen funktionalen Abhängigkeiten der TGA und MSR effizient abzubilden und als Grundlage für ein erweiterbares Inbetriebnahmemanagement zu nutzen?
    \item[FF2:] Wie lässt sich eine modulare Multi-Tenant-Architektur für Inbetriebnahmemanagement-Software gestalten, welche funktionale Erweiterungen ermöglicht?
\end{description}
\subsection{Zielsetzung}
Das primäre Ziel dieser Arbeit ist die Konzeption einer technologischen Architektur,
die den systemimmanenten Vendor-Lock und die Problematik proprietärer Insellösungen durch Transparenz überwindet.
Im Zentrum steht dabei die Nutzung des \ac{IFC}-Standards als universelle Sprache für Bauwerksdaten.
Um die komplexen Beziehungen zwischen Bauteilen, Gewerken und Prozessen performant abzubilden,
wird das System auf einer Graphdatenbank aufgesetzt. Im Gegensatz zu starren relationalen Tabellen
ermöglicht die Graphentechnologie eine flexible Verknüpfung der Datenquellen und bildet so das Rückgrat
für die Interkonnektivität zwischen den Projektbeteiligten.
Ergänzt wird dieser datenzentrierte Ansatz durch eine Add-on-fähige Software-Architektur.
Diese erlaubt es, spezifische Anforderungen für das Inbetriebnahmemanagement als modulare Erweiterungen zu implementieren,
ohne den Software-Kern modifizieren zu müssen. Durch die Kombination aus standardisiertem Datenaustausch (\ac{IFC}) und modularer Erweiterbarkeit
entsteht ein Ökosystem, das nicht nur den Datenaustausch zwischen verschiedenen Gewerken sicherstellt, sondern Unternehmen die
Souveränität über ihre digitalen Prozesse zurückgibt und eine zukunftssichere Skalierbarkeit in Großprojekten ermöglicht.
\subsection{Methodik}
Um die beschriebene Zielsetzung zu erreichen, wird diese Arbeit in drei Teilschritte unterteilt. Zuerst soll durch Literaturrecherche und prototypische Tests ein Domänenmodell mit der Graphdatenbank \textit{Neo4j} und dem Framework \textit{GraphQL} als Vermittlerinstanz
auf seine Vorteile und Machbarkeit untersucht werden. Im nächsten Schritt soll eine modulare Softwarearchitektur, aufbauend auf dem Konzept der Microfrontends, entworfen werden. Die technische Umsetzung soll mit dem Build-Tool \textit{Vite} und dem Architekturmuster
\textit{Module Federation} erfolgen. Das finale Ziel dieses Schrittes ist ein Prototyp als Proof of Concept. Das Zusammenspiel dieser beiden Technologien bietet eine moderne und entwicklerfreundliche Möglichkeit, Microfrontends als Webapplikationen herauszuführen.

Auf diesen beiden Lösungsansätzen aufbauend soll als letzter Schritt die Zusammenführung beider Teilsysteme erfolgen. Hierbei wird der Prototyp der modularen Softwarearchitektur über die GraphQL-Schnittstelle an die Graphdatenbank angebunden.